\title{Byte Latent Transformer: Patches Scale Better Than Tokens}

\begin{abstract}
We present the Byte Latent Transformer ({\textsc BLT}), a novel byte-level large language model architecture that, for the first time, achieves performance on par with tokenization-based LLMs at scale while also delivering notable gains in inference efficiency and resilience. In {\textsc BLT}, bytes are grouped into adaptively sized patches, which act as the central computational entities. Patch segmentation is governed by the predicted entropy of the subsequent byte, allocating greater computational and representational resources in regions of higher data complexity. We provide the inaugural scaling analysis of byte-level models that controls for floating point operations, scaling up to 8B parameters and 4T bytes of training data. Our findings confirm that it is viable to train large models directly on unsegmented byte streams without a predetermined vocabulary. Training and inference speed are enhanced because longer patches are used when data is predictable, resulting in qualitative advances in reasoning and generalization to less common cases. In summary, for a given inference budget, {\textsc BLT} achieves notably improved scaling relative to tokenization-based approaches by expanding both the patch length and model size in tandem.
\end{abstract}

\section{Introduction}
\label{section:intro}

We present the Byte Latent Transformer~(\textbf{BLT}), an architecture that forgoes traditional tokenization and instead directly processes raw byte sequences. BLT is the first model of its kind to achieve parity with tokenization-dependent models at scale, while delivering notable gains in both efficiency and robustness (\S\ref{sec:robustness}). Current large language models (LLMs) are predominantly trained in an end-to-end fashion, with the exception of tokenization—a heuristic, pre-processing mechanism that segments byte streams into a predetermined vocabulary of tokens. This process introduces biases in how information is encoded, resulting in several limitations: susceptibility to input and domain-specific variations~\citep{dagangetting}, vulnerability to noise in the input (\S\ref{sec:robustness}), inadequate handling of orthographic distinctions~\citep{edman2024cute}, and disparities across languages in multilingual settings~\citep{liang2023xlm,petrov2024language,limisiewicz2024myte}.

Tokenization has traditionally played a critical role, as training llms on raw bytes directly incurs unsustainable computational costs at large scales, mainly due to increased sequence lengths~\citep{xue2022byt5}. To address this, earlier research has introduced strategies such as optimizing self-attention mechanisms~\citep{el2020characterbert, clark2022canine} or adopting architectures that eliminate attention entirely~\citep{wang2024mambabyte}~(\S\ref{section:related_work}). Nonetheless, these advances have predominantly benefited the training of \textit{small models}. When scaling up, the primary computational burden of a Transformer arises from the expansive feed-forward network layers—which process every byte—rather than the attention component itself.

To optimize compute usage, we introduce a flexible, learnable approach for aggregating bytes into \textit{patches}~(\S\ref{section:patching}), along with a novel architecture that integrates both byte-level and patch-level features. In contrast to tokenization, {\textsc BLT} does not employ a predetermined patch vocabulary. Instead, any subset of bytes can be transformed into a latent patch embedding through efficient encoder and decoder modules that are trained end-to-end. Our experiments demonstrate that this strategy yields a \textit{more} effective distribution of computational resources compared to models built on tokenization.

Tokenization-based large language models (LLMs) devote equal computational resources to each token, prioritizing uniformity over optimal efficiency. However, this approach assumes that the tokenization scheme aligns with prediction difficulty, which is often not the case since tokens are defined using compression heuristics that may not accurately reflect complexity. The foundational principle of our proposed architecture is to enable the model to selectively apply computational effort according to the demands of each prediction. For instance, predicting the final letters of a word generally constitutes a straightforward, low-entropy task, in contrast to generating the initial word in a sentence, which is typically more complex. Accordingly, {\textsc BLT} introduces a multi-tiered structure~(\S\ref{section:architecture}) comprised of three transformer modules: two compact local models operating at the byte level, and a more substantial global latent transformer~(\autoref{fig:arch_high_level}). 

To facilitate adaptive compute allocation, {\textsc BLT} organizes bytes into patches using the entropy of next-byte predictions. This process yields context-sensitive groupings wherein each patch contains bytes exhibiting similar levels of informational content.

We conduct the inaugural flop-constrained scaling analysis of byte-level models, extending up to 8B parameters and 4T training bytes, and demonstrate that it is feasible to train models directly from bytes on a large scale, thereby avoiding fixed-vocabulary tokenization altogether. Our results indicate that {\textsc BLT} achieves parity with Llama 3 in terms of training performance per flop,\footnote{We measure model computational requirements by the total number of Floating Point OPerations (flops).} but can reduce inference-time flop consumption by as much as 50\% (\S\ref{section:scaling}). Furthermore, our findings reveal that modeling directly on byte sequences yields notable gains for data with infrequent patterns, especially improving the handling of the long-tail. The {\textsc BLT} models demonstrate greater resilience to noisy data relative to their tokenizer-dependent counterparts, and they show improved proficiency at the character level across tasks such as orthography, phonology, and translation in low-resource settings (\S\ref{sec:robustness}). In addition, {\textsc BLT} allows for simultaneous scaling of both the model and the patch size, all while preserving a fixed inference flop constraint. Employing larger patch sizes generally reduces computational demand, thus allowing us to allocate more resources to the global latent transformer by decreasing the frequency of its execution. In inference flop–constrained scaling studies (\autoref{fig:fixed_inference_scaling}), we observe markedly superior scaling behavior compared to models employing tokenization-based architectures.

To summarize, this work offers several key contributions: 1) We present {\textsc BLT}, a byte latent large language model architecture capable of dynamically modulating computational allocation to enhance flop efficiency; 2) We demonstrate that our approach matches Llama 3 in training flop-controlled performance up to the 8B parameter regime, and additionally enables up to 50\% improvements in flop efficiency at the cost of only marginal reductions in evaluation scores; 3) {\textsc BLT} introduces a novel paradigm for scaling large language models, permitting increases in model size without exceeding a predetermined inference compute budget; 4) We provide evidence that {\textsc BLT} models exhibit greater resilience to input noise and display an ability to capture sub-word features in the input that are typically overlooked by token-based large language models. All code for training and inference with {\textsc BLT} is publicly available at~\url{https://github.com/facebookresearch/blt}.

\section{Patching: From Individual Bytes to Groups of Bytes}
\label{section:patching}

Dividing bytes into discrete \textit{patches} enables {\textsc BLT} to flexibly distribute computational resources according to contextual needs. As illustrated in Figure~\ref{fig:patching_types}, there exist multiple approaches for organizing bytes into patches. More precisely, a patching function $f_p$ partitions a byte sequence $\pmb{x}=\{x_i,|i=1,\ldots n\}$ of length $n$ into a set of $m < n$ patches $\pmb{p}=\{p_j|j=1,\ldots,m\}$ by assigning each $x_i$ to either 0 or 1, where a 1 marks the commencement of a new patch. Whether the model is token-based or patch-based, the primary determinant of computational demand is the number of steps performed by the main Transformer module. Within {\textsc BLT}, this corresponds to the count of patches produced by the selected patching function for representing the data. As a result, the mean patch length, referred to as \textit{patch size}, plays a crucial role in estimating the processing cost during training and inference for a specific patching scheme~(\S\ref{section:flops}).

Subsequently, we outline three distinct patching strategies: assigning a fixed byte count per patch~(\S\ref{section:static-patch}), grouping by whitespace boundaries~(\S\ref{section:white-patch}), and adaptively defining patches using entropy metrics derived from a compact byte-level language model~(\S\ref{section:dyn-patch}). In closing, we explore incremental patching mechanisms and clarify distinctions between patching and traditional tokenization methods~(\S\ref{section:bpe-patch}).

\subsection{Strided Patching Every K Bytes}
\label{section:static-patch}

One of the simplest strategies for aggregating bytes involves dividing them into fixed-size patches of length $k$, as employed in MegaByte~\citep{yu2023megabyte}. This approach, based on a constant stride, offers ease of implementation during both training and inference phases, provides a direct method to adjust the mean patch size, and thus allows straightforward regulation of computational complexity. Nevertheless, this fixed patching technique introduces notable limitations. Firstly, computational resources are distributed uniformly rather than adaptively, resulting in potential inefficiency—such as squandering transformer step $j$ on regions consisting solely of whitespace in code, or failing to devote enough processing power to highly informative segments like those containing mathematical content. Secondly, this method produces variability and a lack of contextual consistency in how similar byte sequences are partitioned; for instance, the same word might be segmented in dissimilar ways.

\subsection{Space Patching}
\label{section:white-patch}

According to \citet{slagle2024spacebyte}, a straightforward but effective refinement to strided patching is introduced, wherein new patches are generated immediately following any space-like bytes\footnote{
    Space-like bytes encompass all bytes except latin letters, digits, or utf-8 continuation bytes; moreover, each patch must include at least one non space-like byte.}, which tend to correspond to natural linguistic unit boundaries in numerous languages. The Space patching approach devotes a latent transformer operation (i.e., higher computational cost) to each word, thus assigning increased modeling capacity per word. This method guarantees consistent patching for words throughout sequences, and dedicates more computational effort to predictions at difficult points—typically those appearing after a space. For instance, determining the first byte of a response to the prompt ``Who composed the Magic Flute?\uline{\hspace{.6em}}'' involves much greater uncertainty than predicting the subsequent bytes after the ``M'' in ``Mozart,'' since the initial character significantly narrows the set of plausible options, rendering the remainder of the completion easier to anticipate. Nonetheless, space patching exhibits limitations: it does not adapt smoothly to all languages or domains, nor can it adjust the size of its patches. In the subsequent section, we present a novel patching approach inspired by the observation that initial bytes of words are generally the hardest to predict, which also enables straightforward control over patch size.

\subsection{Entropy Patching: Using Next-Byte Entropies from a Small Byte LM}
\label{section:dyn-patch}

Instead of depending on heuristic methods like whitespace detection, we adopt a data-centric strategy to determine areas of high uncertainty in next-byte predictions. Our proposed technique, \textit{entropy patching}, leverages entropy measurements to establish the locations of patch boundaries.

To this end, we fit a compact byte-level auto-regressive language model on the training corpus for {\textsc BLT}, enabling us to calculate the entropy of predicting the subsequent byte according to the model’s distribution $p_{e}$ over the byte-level vocabulary $\mathcal{V}$:

\begin{equation}
H(x_i) = \sum_{v \in \mathcal{V}} p_{e}(x_i=v|\pmb{x}_{<i}) \log p_{e}(x_i=v|\pmb{x}_{<i})
\end{equation}

We employ two techniques to detect patch boundaries using the entropies $H(x_i)$. The first method selects points that exceed a predetermined global entropy threshold, as shown in~\autoref{fig:patching}. The second method targets points whose entropy is significantly higher compared to that of the preceding point. This latter strategy can also be understood as marking locations where the roughly monotonic decrease in entropy within a patch is disrupted.

\begin{align*}
    \text{Global Constraint}\hspace{1cm}H(x_t)& > \theta_g\\
    \text{Approx. Monotonic Constraint}\hspace{1cm}H(x_t)& - H(x_{t-1}) > \theta_r
\end{align*}

Patch boundaries are determined through a simple preprocessing process performed at the time of data loading. This approach contrasts with~\citet{nawrot-etal-2023-efficient}, where a classifier is utilized to identify entropy-based patch boundaries. In our study~(\S\ref{section:experiments}), we evaluate and contrast these two approaches for separating bytes characterized by low and high entropy.

\subsection{The Byte-Pair Encoding (BPE) Tokenizer and Incremental Patching}
\label{section:incremental-patching}
\label{section:bpe-patch}

A significant number of contemporary llms, such as our baseline Llama 3, utilize subword tokenizers like bpe~\citep{gage1994new, sennrich-etal-2016-neural}. In our terminology, ``tokens'' are defined as groups of bytes selected from a \textit{finite} vocabulary established before model training, contrasting with ``patches,'' which are constructed dynamically from byte sequences and do not rely on a fixed vocabulary. An essential distinction is that when using tokens, the model is deprived of direct access to the raw byte-level information.

An important advancement introduced by {\textsc BLT} relative to tokenization-based approaches lies in its reconfiguration of the balance between vocabulary size and computational requirements. Traditional llms face a trade off: expanding the vocabulary leads to larger average token sizes, which can reduce the total number of steps required for processing. However, this also results in a substantial increase in the output dimensionality of the model's final projection layer. As a consequence, tokenization-based models are limited in their ability to significantly alter token sizes without incurring prohibitive increases in inference costs. As an illustration, Llama 3 achieves a rise in average token size from 3.7 to 4.4 bytes, but this comes at the expense of quadrupling the size of its embedding table relative to Llama 2.

During the generation process, {\textsc BLT} must determine at each point in the byte sequence whether the current position marks a patch boundary, since this influences whether the Latent Transformer needs to be activated for further computation. Importantly, this assessment must be made without reference to subsequent, as-yet-unproduced bytes; in other words, patching decisions cannot rely on future context to define the segmentation of the sequence. More formally, a patching method $f_p$ exhibits the property of incremental patching if it adheres to the following condition:
$$f_p(\pmb{x}_{<i}) = f_p(\pmb{x})_{<i}$$
bpe does not fulfill the requirements of an incremental patching scheme, because the same leading subsequence may be tokenized differently based on the subsequent bytes, and thus the property above is violated\footnote{Introducing a dedicated delimiter token to indicate patch boundaries can make bpe incremental, but this comes at the cost of lengthening the byte sequence.}.

\section{{\textsc BLT} Architecture}
\label{section:architecture}
{\textsc BLT} consists of a substantial global autoregressive language model that processes representations at the patch level, and is complemented by two more compact local modules responsible for encoding byte sequences into patches and for reconstructing bytes from patch representations~(\autoref{fig:arch_high_level}).

\subsection{Latent Global Transformer Model}

\textit{The Latent Global Transformer} refers to an autoregressive transformer, denoted as $\mathcal{G}$ and comprised of $l_{\mathcal{G}}$ layers, which transforms a series of latent patch representations, $p_j$, into corresponding output patch representations, $o_j$. In this work, the index $j$ is consistently reserved for patches, while $i$ refers specifically to bytes. The global transformer implements a block-causal attention mask~\citep{dubey2024llama}, ensuring that each position can attend only to itself and previous patches within the same document. The majority of computational overhead, during both pre-training and inference, is attributed to this model. Consequently, selectively activating the global transformer enables dynamic adjustment of computational resources devoted to distinct input and output segments, tailored to their respective complexities.

\subsection{Local Encoder}
\label{section:local}

\textit{The Local Encoder Model}, represented as $\mathcal{E}$, employs a transformer framework characterized by a relatively small number of layers, specifically $l_{\mathcal{E}} << l_{\mathcal{G}}$. Its central objective is to rapidly convert input byte sequences $b_i$ into dense and informative patch vectors, $p_j$. A notable modification to the standard transformer setup is the introduction of a cross-attention mechanism subsequent to each transformer block, designed to aggregate byte-level features into patch-level representations~(\autoref{fig:crossattn}).

Initially, bytes $b_i$ are projected into continuous vectors using an embedding table of size $\mathbb{R}^{256 \times h_{\mathcal{E}}}$, producing embeddings $x_i$. These vectors can be further enhanced with supplemental information through hash-embeddings, as discussed in~(\S\ref{sec:hash-emb}). The embedded bytes are subsequently passed through a sequence of alternating transformer layers and cross-attention modules~(\S\ref{sec:enc-cross-attn}), which iteratively refine the representations to yield the final patch vectors, $p_i$, utilized as input to the global transformer, $\mathcal{G}$.

Within each transformer block, the attention mechanism enforces a \textit{local block causal} mask: each byte is allowed to focus on up to $w_{\mathcal{E}}$ prior bytes, irrespective of dynamic patch segmentation, but is restricted from attending to bytes beyond the confines of its respective document. The subsections that follow elaborate on the structure of the embeddings and the design of the cross-attention module.

\subsubsection{Encoder Hash n-gram Embeddings}
\label{sec:ngram-emb}
\label{sec:hash-emb}
To produce rich and resilient representations at every position $i$, it is essential to encode the context provided by previous bytes. In {\textsc BLT}, this is accomplished by simultaneously capturing the identity of byte $b_i$ in isolation as well as its role within a byte n-gram. Specifically, at each step $i$, we begin by forming byte-grams

\begin{align}
    g_{i,n}=\{b_{i-n + 1},\ldots, b_i\}
\end{align}

for each position $i$ within the byte sequence and for values of $n$ ranging from three up to eight.\footnote{We do not include $n$-grams of length $n$ or greater when $i<n$.}

Subsequently, we propose hash-based $n$-gram embeddings, in which every byte $n$-gram is assigned an index within a fixed-size embedding table $E_{n}^{hash}$ through a hash function, for each $n \in\{3,4,5,6,7,8\}$~\citep{bai2010learning}. The embedding retrieved in this manner is combined with the individual byte embedding; this sum is then normalized and provided as input to the local encoder architecture. The enriched embedding is computed as

\begin{align}
e_i &= x_i + \sum_{n = 3,...,8}E_{n}^{hash}(\text{Hash}(g_{i,n})) \\
\text{where, Hash}(g_{i,n}) &= \text{RollPolyHash}(g_{i,n}) \% |E_{n}^{hash}|
\end{align}

The value $e_i$ is scaled by dividing by the total count of $n$-gram sizes, incremented by one, and we utilize RollPolyHash according to the definition provided in Appendix~\ref{appendix:rollpolyhash}. Section~\ref{section:ablations} examines the influence of $n$-gram hash embeddings, exploring various choices for $n$ as well as different embedding table sizes, on the observed scaling law patterns under fixed FLOP budgets. Beyond employing hashed $n$-gram embeddings, we have also conducted experiments with frequency-based $n$-gram embeddings; further specifics regarding these experiments are included in Appendix~\ref{appendix:ngrams}.

\subsubsection{Encoder Multi-Headed Cross-Attention}
\label{sec:enc-cross-attn}
Our approach adheres to the structure of the input cross-attention mechanism as described in the Perceiver architecture~\citep{jaegle2021perceiver}, with a notable distinction: instead of using a constant set of latent vectors, the latent representations in our design are dynamic, representing variable patch representations~(\autoref{fig:crossattn}), and each focuses exclusively on the bytes contained within its specific patch. The construction of this module starts with creating a query vector for each patch $p_j$. This query is formed by applying a pooling operation to the byte representations of the patch $p_j$ and subsequently projecting it through a linear transformation, $\mathcal{E}_{C} \in \mathbb{R}^{h_{\mathcal{E}} \times (h_{\mathcal{E}} \times U_{\mathcal{E}})}$, where $U_{\mathcal{E}}$ is the number of heads in the encoder's cross-attention. Letting $f_{\text{bytes}}(p_j)$ represent the sequence of bytes associated with the patch $p_j$, the computation proceeds as follows:

\begin{align}
P_{0,j} &= \mathcal{E}_{C}(f_\text{bytes}((p_j)), f~ \text{is a pooling function} \\
P_l &= P_{l-1} + W_o\left(\text{softmax}\left(\frac{QK^T}{\sqrt{d_k}}\right)V\right) \\
\text{where } Q_j &= W_q(P_{l-1,j}), K_i = W_k(h_{l-1,i}), V_i = W_v(h_{l-1,i}) \\
h_l &= \text{Encoder-Transformer-Layer}_l(h_{l-1})
\end{align}

In this context, $P \in \mathbb{R}^{n_p \times h_{\mathcal{G}}}$ denotes the representations for $n_p$ patches that will be fed into the global model; $P$ is constructed by aggregating the byte-level embeddings $e_i$ for each patch $p_j$. The matrices $W_q$, $W_k$, $W_v$, and $W_o$ serve as linear transformation weights for generating queries, keys, values, and output projections, respectively, with keys and values derived from the byte representations $h_i$ output by the preceding layer (using $e_i$ for the initial layer). To ensure that information from different patches remains distinct, a masking mechanism is applied such that each query $Q_j$ interacts solely with keys and values corresponding to bytes within the same patch $j$. Given that multi-headed attention is applied over $Q, K,$ and $V$, and considering patch representations typically possess higher dimensionality ($h_{\mathcal{G}}$) than the individual embedding size $h_{\mathcal{E}}$, we preserve $P_l$ as several attention heads with dimension $h_{\mathcal{E}}$ during cross-attention, then concatenate these head outputs to form vectors in $h_{\mathcal{G}}$-dimensional space. Furthermore, we incorporate pre-LayerNorm normalization on queries, keys, and values, and intentionally omit positional embeddings in the cross-attention mechanism. The cross-attention block also employs a residual connection around its operations.

\subsection{Local Decoder} 

Analogous to the local encoder, the local decoder $\mathcal{D}$ is also implemented as a lightweight transformer architecture, comprising $l_{\mathcal{D}}$ layers where $l_{\mathcal{D}} << l_{\mathcal{G}}$. This component is responsible for converting the global patch representations $o_j$ into raw bytes, denoted $y_i$. The local decoder generates the sequence of bytes by conditioning on the sequence of previously reconstructed bytes, utilizing the hidden states derived from the local encoder for these sequences. Its structure alternates between $l_{\mathcal{D}}$ cross-attention layers and transformer layers. In each block, the cross-attention module precedes the transformer layer, initially mapping the patch representations to byte-level representations, after which the transformer module processes the constructed byte sequence.

\subsubsection{Decoder Multi-headed Cross-Attention} 

Within the decoder’s cross-attention mechanism, the byte representations function as queries, while the patch representations serve as keys and values, effectively reversing their roles compared to the encoder. To initiate the cross-attention process, byte representations are initialized using the byte embeddings output by the final encoder layer, denoted as $h_{l_{\mathcal{E}}}$. For each subsequent layer $l$, the byte representations, $d_{l,i}$, are updated as follows:

\begin{align}
D_0 &= h_{l_{\mathcal{E}}} \\
B_l &= D_{l-1} + W_o\left(\text{softmax}\left(\frac{QK^T}{\sqrt{d_k}}\right)V\right), \\
  &\text{where } Q_i = W_q(d_{l-1,i}), K_i = W_k(\mathcal{D}_C(o_j)), V_i = W_v(\mathcal{D}_C(o_j)) \\
  D_l &= \text{Decoder-Transformer-layer}_l(B_l)
\end{align}

Here, as before, $W_k$ and $W_v$ denote the key and value projection matrices, respectively, which perform a linear mapping followed by the split operation $\mathcal{D}_C$ on the final patch representations $o_j$ produced by the global model. The query projection matrix $W_q$ is applied to the byte-level representations $d_{l-1}$ from the prior decoder transformer layer (or $h_{l_{\mathcal{E}}}$ for the initial layer). The resulting outputs are then projected through $W_o$. Consequently, $B$ lies in $\mathbb{R}^{h_{\mathcal{D}} \times n_b}$, where $n_b$ represents the total number of output bytes. The subsequent decoder layer representations, $D_l$, are obtained by applying a decoder transformer layer to $B$, the result of the cross-attention stage. Similar to the approach taken in the local encoder’s cross-attention, this architecture leverages multiple attention heads, incorporates pre-layer normalization, omits positional embeddings, and includes a residual pathway around the cross-attention block.

\section{Experimental Setup}
\label{section:experiments}
We establish a series of rigorously controlled experiments to evaluate {\textsc BLT} in relation to models that utilize tokenization, ensuring that {\textsc BLT} does not receive any unfair benefit from potential increases in effective sequence length.

\subsection{Pre-training Datasets} The models evaluated in this study are all initially trained on two principal datasets: (1) The Llama 2 dataset~\citep{touvron2023llama}, containing 2 trillion tokens sourced from a range of publicly accessible materials, which are rigorously cleaned and filtered to ensure higher data quality; and (2) {\textsc BLT}-1T: An original corpus assembled for this work, comprising 1 trillion tokens collected from diverse public domains and including a portion of the pre-training set made available by Datacomp-LM~\citep{li2024datacomp}. The Llama 2 dataset serves primarily in our scaling law investigations—specifically for experiments evaluating the optimal number of training tokens as per~\cite{dubey2024llama}—to guide architectural decisions for {\textsc BLT}. In contrast, {\textsc BLT}-1T is employed for end-to-end pre-training to facilitate comparison with Llama 3 on downstream benchmarks. Both datasets exclude any information originating from Meta products or services. Additionally, for our tokenizer-based baselines, we utilize the Llama 3 tokenizer featuring a 128K token vocabulary, as it demonstrated superior baseline results in our trials compared to the Llama 2 tokenizer.

\subsection{Entropy Model} For our experiments, the entropy model consists of a byte-level language model, which is trained using the same data distribution as the full {\textsc BLT} model. By default, the configuration includes a transformer architecture with 100 million parameters, composed of 14 layers and a hidden size of 512. The attention mechanism uses a sliding window of 512 bytes. All other hyperparameter choices align with those used in both the local and global transformer models. As elaborated in \autoref{section:ablations}, we explored variations in model capacity, receptive field size, and model architecture. Notably, when the receptive field is sufficiently constrained, the entropy model can be represented compactly via a lookup table.

\subsection{Entropy Threshold and Equalizing Context Length} For models employing entropy-based patching, we determine a patching threshold that yields a specified mean \textit{patch size} over the pretraining data mixture. In {\textsc BLT}, as opposed to tokenization methods, the \textit{patch size} can be set to any value, which considerably impacts the context window utilized by the model. To keep the average context length consistent and prevent models with larger patch sizes from benefiting disproportionately, we ensure that the expected number of bytes per batch stays fixed. As a result, for models with increased patch sizes, the input sequence length is reduced accordingly. Specifically, on Llama 2 data, a context window of 8k bytes is employed, while for the {\textsc BLT}-1T dataset, the context length is expanded to an average of 16k bytes, yet the batch size remains at a constant expected value of 16M bytes.

Although the mean batch size remains unchanged, dynamic patching approaches result in varying proportions of bytes relative to patches within each batch. To improve computational efficiency, our {\textsc BLT} training framework assembles batches such that each contains a uniform number of patches, thereby eliminating the need for padding during the more resource-intensive latent transformer step. To prevent excessive memory usage caused by exceptionally large patches, we apply padding or truncation to input byte sequences: for Llama 2 datasets, sequences are standardized to 12k bytes, while for {\textsc BLT}-1T datasets, the limit is set to 24k bytes.

\subsection{Entropy Model Context}
\label{section:entropy-context}
Our empirical observations indicate that entropy patching tends to result in increasingly larger patches when applied to structured formats such as multiple choice tasks, which are frequently characterized by repetitive patterns (see patching on an MMLU instance in \autoref{fig:mmlu_patching}). These differences arise due to reduced entropy within sections of content that repeat in the entropy model context. Therefore, during the extensive evaluation of {\textsc BLT}-Entropy with a patch size of 4.5, we refresh the entropy context by inserting new lines and employ an approximate monotonicity constraint, since it is less susceptible to "entropy drift" induced by fluctuations in context length. Notably, this adjustment impacts only the methodology for entropy calculation; our approach for selecting the entropy threshold remains consistent.

\subsection{FLOPs Estimation} 
\label{section:flops}

Our approach for calculating transformer flops is primarily based on the methodology described in Chinchilla~\citep{hoffmann2022training}, which accounts for the flops arising from feed-forward networks, the qkvo projections within the self-attention mechanism, as well as the attention calculation and output projection steps. However, unlike the original framework, we treat the input embedding layer as an efficient lookup rather than a dense matrix multiplication, effectively assigning it zero flops. Consistent with established literature, we assume that the backward propagation step incurs double the flops of the forward propagation.

To determine the flops \textit{per byte} for {\textsc BLT} architectures, we aggregate the floating-point operations involved in the local encoder transformer, global latent transformer, and local decoder transformer components, including the cross attention mechanisms within both the encoder and decoder stages:

\begin{align}
    \text{FL}_{\text{{\textsc BLT}}} &= \text{Transf. FL}(h_{\mathcal{G}}, l_{\mathcal{G}}, m=n_{ctx}/n_p,V=0) / n_p\\
   &+\text{Transf. FL}(h_{\mathcal{E}}, l_{\mathcal{E}}, m=w_{\mathcal{E}}, V=0) \\
   &+ \text{Transf. FL}(h_{\mathcal{D}}, l_{\mathcal{D}}, m=w_{\mathcal{D}}, V=256) \\ 
    &+ \text{Cross Attn. FL}(h_{\mathcal{E}}, l_{\mathcal{E}}, m=n_p, r=n_p/k) \times k/n_p \\ 
   &+ \text{Cross Attn. FL}(h_{\mathcal{D}}, l_{\mathcal{D}}, m=k, r=k/n_p) 
\end{align}

Here, $n_{ctx}$ denotes the sequence length measured in bytes, while $n_p$ refers to the patch size. The parameter $r$ represents the proportion of queries relative to keys/values. The variable $k$ signifies the ratio between patch-dimension and byte-dimension, that is, it reflects how many local model splits are merged to compose a single global representation ($k=2$ in \autoref{fig:crossattn}). $V$ stands for the vocabulary size employed by the output projection in the local decoder exclusively. The context length $m$ used in the attention mechanism varies based on whether the module processes the byte or patch sequence. We update the attention FLOPs accordingly for every component. Detailed formulae for calculating Transformer-FLOPs and Cross-Attention FLOPs are presented in Appendix~\ref{section:flops-appendix}.

\subsection{Bits-Per-Byte Estimation} Perplexity serves as a meaningful metric only when a specific tokenizer is employed, since it quantifies uncertainty at the token level. To facilitate comparisons between byte-level and token-level language models, as suggested in prior research~\citep{xue2022byt5, yu2023megabyte, wang2024mambabyte}, we utilize Bits-Per-Byte (BPB) as our evaluation metric. BPB provides a measure analogous to perplexity, but independent of any particular tokenization scheme. The formal definition is as follows:

\begin{align}
    \text{BPB}(x)&= \frac{\mathcal{L}_{CE}(\pmb{x})}{\ln(2)\cdot n_{\text{bytes}}}
\end{align}

Here, the aggregate cross-entropy loss computed over the sequence $\pmb{x}$ represents model uncertainty, which is then standardized by dividing by both the total number of bytes in $\pmb{x}$ and the normalization constant.

\subsection{Transformer Architecture Hyperparameters}  
For the transformer modules in {\textsc BLT}, encompassing both local and global variants, we primarily adopt the configuration employed in Llama~3~\citep{dubey2024llama}. Specifically, the feed-forward networks utilize the SwiGLU activation~\citep{swiglu}, while self-attention layers incorporate rotary positional embeddings (RoPE)~\citep{RoPe} with a parameter setting of $\theta=500000$~\citep{xiong2024effective}. Layer normalization is carried out via RMSNorm~\citep{RMSNorm}. All self-attention modules that rely on static attention masks, such as \textit{block causal} or \textit{fixed-window block causal} masks, leverage Flash attention~\citep{dao2022flashattention}; in these scenarios, we set the attention window size to 512 for fixed-width masking. For cross-attention layers—where the masking is dynamically determined based on patches—we employ Flex Attention\footnote{\url{https://pytorch.org/blog/flexattention}}, which enables the use of optimized fused kernels to substantially accelerate training.

\subsection{{\textsc BLT}-Specific Hyperparameters}

To evaluate the performance of {\textsc BLT} models, we perform experiments focusing on two main aspects: scaling behaviors and performance on downstream tasks. Our analysis encompasses model variants with parameter counts of 400M, 1B, 2B, 4B, and 8B. The corresponding architecture-related hyperparameters for each configuration can be found in Appendix~Table~\ref{tab:arch}.

For initialization in the first cross-attention layer of the local encoder, we employ max-pooling for the query vectors. All {\textsc BLT} models use $500,000$ hashes with a single hash function and n-gram lengths varying from 3 to 8. The learning rate across all models is set at $4e-4$. This value was selected based on a hyperparameter sweep over learning rates between $1e-3$ and $1e-4$ at the 400M and 1B parameter scales, which revealed that the optimal setting was the same for both token-based and {\textsc BLT} architectures. Training batch sizes on Llama-2 data adhere to the recommendations of~\cite{dubey2024llama}, or the equivalent when measured in bytes. 

For the optimization process, AdamW~\citep{adamw} is utilized, configured with $\beta_1$ of 0.9, $\beta_2$ of 0.95, and $\epsilon=10^{-8}$. The training schedule includes a linear warm-up over 2000 steps, followed by a cosine decay of the learning rate to zero. Weight decay is maintained at 0.1, and a global gradient clipping threshold of 1.0 is applied.

\section{Scaling Trends}
\label{section:scaling}

We offer a comprehensive overview of scaling patterns in byte-level models, providing insights that can guide the future expansion of {\textsc BLT} models. Our investigation into scaling seeks to overcome shortcomings of earlier studies on byte-level approaches by: (a) examining scaling behaviors within compute-optimal training frameworks, (b) training equivalent 8B parameter models with substantial datasets—reaching up to 1T tokens (or 4T bytes)—and assessing performance on a variety of downstream tasks, and (c) analyzing scaling behaviors under constraints where inference costs are held constant. Subsequent sections will explore the distinct benefits that arise from processing byte-sequences directly.

\subsection{Parameter Matched Compute Optimal Scaling Trends}
\label{section:parameter-matched-scaling}
Utilizing the Llama 2 dataset, we train a set of \textit{compute-optimal} bpe and {\textsc BLT} models at four parameter scales, spanning from 1B to 8B parameters. Training flops are then compared with language modeling accuracy on a representative subset drawn from the training data mixture. The bpe models are configured according to the optimal ratio of parameters to training data, following guidelines established by Llama~3~\citep{dubey2024llama}. This \textit{compute-optimal} configuration is theoretically expected to deliver maximal performance on the training data within a specified computation budget~\citep{hoffmann2022training}, thereby serving as a strong point of reference for our analysis. In addition, for every bpe model, we construct a matched {\textsc BLT} variant, employing a Latent Transformer identical in parameter count and architecture to its corresponding bpe Transformer, and train it on the same dataset.

As shown in~\autoref{fig:scaling-compute-opt} (right), {\textsc BLT} models consistently equal or surpass the performance of their bpe-based equivalents, and this relationship persists as model size and computational resources increase. According to our current understanding, {\textsc BLT} represents the first byte-level Transformer model to exhibit scaling behaviors comparable to BPE-based architectures when operating under compute-optimal conditions. These findings lend support to our hypothesis that the ideal parameter-to-training-compute ratio identified for bpe models is also applicable to {\textsc BLT}, or at minimum, does not differ significantly.

Advancements in both architecture and dynamic patching are essential for achieving scaling patterns that align with bpe models. In ~\autoref{fig:scaling-compute-opt} (left), we present a comparison between models that utilize space-patching techniques and Llama 3. For this analysis, we approximate SpaceByte \citep{slagle2024spacebyte} as a {\textsc BLT} space-patching variant, omitting n-gram embeddings and cross-attention modules. While SpaceByte demonstrates gains relative to Megabyte, there remains a notable gap compared to Llama 3's performance. The panel on the right in \autoref{fig:scaling-compute-opt} demonstrates the cumulative benefits arising from both structural innovations and the implementation of dynamic patching. At scale, {\textsc BLT} models exhibit performance comparable to leading tokenizer-based approaches such as Llama 3.

Additionally, we find that the specific choice of tokenizer influences model outcomes: models trained using the Llama-3 tokenizer achieve superior results relative to those employing the Llama-2 tokenizer on identical training data.

In summary, the {\textsc BLT} architecture demonstrates an intermediate scaling behavior between Llama 2 and 3 when employing much larger patch sizes. While the bpe tokenizers of Llama 2 and 3 exhibit average token sizes of 3.7 and 4.4 bytes respectively, {\textsc BLT} is able to maintain similar scaling dynamics even with average patch sizes of 6 or 8 bytes. Since the number of inference flops decreases as patch size increases, adopting an 8-byte patch size can result in nearly 50\% reduction in inference flops. Furthermore, models utilizing larger patch sizes tend to show improved performance as both model and dataset sizes increase. Specifically, {\textsc BLT} with an 8-byte patch size begins with notably lower performance than the bpe Llama 2 model at the 1B parameter scale, but surpasses bpe Llama 2 when scaled up to 7B parameters. This observation indicates that larger patch sizes may yield superior results at greater model scales, and potentially even larger patch sizes could be practical as model size and training computation continue to expand.

\subsection{Beyond Compute Optimal Task Evaluations}
\label{section:evals}
To further investigate scaling characteristics, we train an 8B {\textsc BLT} model past the compute-optimal threshold using the {\textsc BLT}-1T dataset, which comprises a larger volume of higher-quality data. We then evaluate the model’s capabilities across a diverse array of widely recognized classification and generation benchmarks. For this assessment, we focus on tasks that examine common sense reasoning, general world knowledge, and code generation abilities:
\paragraph{Classification tasks} encompass ARC-Easy (0-shot)~\citep{Eval_ARC}, Arc-Challenge (0-shot)~\citep{Eval_ARC}, HellaSwag (0-shot)~\citep{Eval_hellaswag}, PIQA (0-shot)~\citep{Eval_piqa}, and MMLU (5-shot)~\citep{hendrycks2020measuring}. We utilize a prompt-scoring strategy based on the model’s likelihood estimates over the answer options and present mean accuracy results.
\paragraph{Coding related generation tasks:} To measure the proficiency of LLMs in Python code synthesis, we provide pass@1 metrics for MBPP (3-shot)~\citep{Eval_MBPP} and HumanEval (0-shot)~\citep{Eval_HumanEval}.

In \autoref{tab:evals}, the performance of three models trained on the {\textsc BLT}-1T dataset is presented: a model based on the Llama 3 tokenizer using BPE,\footnote{The Llama 3 tokenizer, with its 128k vocabulary, was selected for its superior performance compared to the 32k vocabulary of Llama 2.} along with two {\textsc BLT} model variants. The first variant applies a space-patching strategy ({\textsc BLT}-Space), while the second leverages an entropy-based patching approach ({\textsc BLT}-Entropy), both incorporating an approximate monotonicity constraint, and for the entropy-based model, context resets are triggered by new lines as elaborated in~\autoref{section:entropy-context}. All models are trained using the same flop budget. For {\textsc BLT}-Entropy, we further lower the entropy threshold at inference from 0.6 to 0.1, which enhances task performance but incurs additional inference steps.

The {\textsc BLT}-Entropy model achieves superior results compared to the Llama 3 model in 4 of the 7 evaluated tasks, despite both models being trained with an equivalent amount of data in bytes. This performance gain likely stems from two main factors: (1) more efficient utilization of computational resources during training through dynamic patching, and (2) the explicit modeling of data at the byte level rather than at the token level.

Conversely, {\textsc BLT}-Space lags behind the Llama 3 tokenizer in performance across all but a single task, yet it notably lowers inference flops due to its increased mean patch size of 6 bytes. By contrast, both the bpe and entropy-patching based models exhibit a similar average patch size of about 4.5 bytes within the training dataset mixture. Given an equivalent training budget, the model employing larger patches processes 30\% more data than the other two, potentially moving {\textsc BLT} further from achieving compute-optimality.

\subsection{Patches Scale Better Than Tokens} The {\textsc BLT} architecture enables the joint scaling of both model and patch sizes without altering the total flop budget for training and inference, nor requiring more training data. Patch-based approaches, unlike fixed-vocabulary token-based models, afford the flexibility to enlarge patch size arbitrarily, circumventing the inherent efficiency compromises described in Section~\ref{section:bpe-patch}. By increasing patch length, computational demand is reduced, which, in turn, permits additional resources to be redirected toward enlarging the global latent transformer, as it is invoked less frequently.

We perform a fixed inference scaling analysis to investigate whether increasing model size and using fewer inference steps on larger patches yields superior performance compared to utilizing smaller models with more steps. Beginning with models containing 400 million and 3.6 billion parameters, both utilizing the Llama 2 tokenizer, we identify flop-equivalent models that employ the Llama 3 tokenizer as well as {\textsc BLT}-Entropy architectures featuring average patch sizes of 6 and 8 bytes on the training data mixture (refer to~\autoref{tab:fixed-inf-params} for a detailed breakdown of the models). For models with patch size 8, three encoder layers are employed instead of a single layer. We subject each configuration to training across multiple predetermined flop budgets.

As depicted in \autoref{fig:fixed_inference_scaling}, {\textsc BLT} models demonstrate superior scaling behavior compared to tokenization-based systems across both classes of inference flops. Initially, BPE models yield stronger results when training resources are limited; however, as training budgets increase beyond the compute-optimal threshold, {\textsc BLT} models quickly surpass them in performance. In real-world applications, it is often advantageous to allocate greater resources to the one-time pretraining phase in order to obtain models that perform better under a fixed inference cost. This approach is exemplified by 8B model variants such as Llama 3.1, which has undergone training on data volumes that are two orders of magnitude larger than the compute-optimal quantity for its parameter count.

The point at which {\textsc BLT} surpasses token-based architectures has shifted somewhat nearer to the compute-optimal regime with larger model scales—changing from 3x to 2.5x above the compute-optimal budget for the highest flop class considered. Additionally, the model employing a larger patch size of 8 demonstrates more pronounced scaling in the highest flop class and exceeds the performance of smaller patch size models at an earlier stage. As elaborated in Section~\ref{section:parameter-matched-scaling}, increasing patch sizes tend to narrow the performance gap with BPE models when evaluated at larger model sizes. This phenomenon appears to result, at least in part, from the byte-level Encoder and Decoder components constituting a decreasing proportion of overall compute, as their computational requirements scale more slowly compared to the Latent Transformer. For example, when scaling up the total parameter count by 20x (from 400M to 8B), the parameters dedicated specifically to {\textsc BLT}'s local modules only about double in number. This distinction is significant because larger patch sizes impact only the flops associated with the patch Latent Transformer, rather than those from byte-level components. Indeed, this helps explain why the parameter ratio for {\textsc BLT}-Entropy with ps=8 increased from 1.6x to 1.7x relative to the Llama 2 model size at greater model scales.

In conclusion, our investigation into patch-length scaling reveals that the {\textsc BLT} patch-based architecture benefits from joint increases in both patch size and model size, resulting in more favorable scaling behavior. These improvements are not only sustained but may also be enhanced as the model scale grows.

\section{Byte Modeling Improves Robustness} Furthermore, we evaluate the robustness of {\textsc BLT} in comparison with token-based architectures that do not utilize byte-level features, and introduce a technique to augment pretrained token-based models with byte-level capabilities.
\label{sec:robustness}

\subsection{Character-Level Tasks}

One of the initial incentives for developing byte-level models was their inherent resilience to noise present at the byte level in the input, in addition to their ability to recognize sub-token structures—a capability often lacking in models that rely on traditional tokenizers. To systematically assess these properties, we conduct further experiments on benchmarks specifically designed to evaluate both noise robustness and sensitivity to character-level information, covering English and multiple languages, and incorporating elements such as digits and phonemes. The outcomes of these assessments are detailed in Table~\ref{tab:char_tasks}.

\paragraph{Noisy Data} To assess the resilience of different modeling approaches, we introduce artificially corrupted variants of the benchmark classification datasets outlined in Section~\ref{section:evals}. Our comparison focuses on tokenizer-based models and {\textsc BLT}. Five character-level noise techniques are utilized to perturb the input text: (a) \textit{AntSpeak}, which transforms all text into uppercase letters separated by spaces; (b) \textit{Drop}, where 10\% of the characters are randomly deleted; (c) \textit{RandomCase}, randomly alternating the casing of 50\% of the characters to uppercase and 50\% to lowercase; (d) \textit{Repeat}, which duplicates 20\% of the characters up to four consecutive times; and (e) \textit{UpperCase}, which changes all characters in the text to uppercase. During evaluation, we systematically apply each noise method to the prompt, the completion, or both, treating each scenario as a separate task and averaging the results. As shown in Table \ref{tab:char_tasks}, evaluations on noised HellaSwag~\citep{Eval_hellaswag} demonstrate that {\textsc BLT} consistently shows superior robustness over tokenizer-based models, outperforming a comparable baseline by an average of 8 points, and even surpassing the Llama 3.1 model, which was trained on a much more extensive dataset.

\paragraph{Phonology - Grapheme-to-Phoneme (G2P)} We evaluate {\textsc BLT}'s proficiency in converting sequences of graphemes (alphabetic representations of words) into their corresponding phonemic transcriptions (representations of pronunciation). Table \ref{tab:char_tasks} displays the outcomes of the G2P task under a 5-shot configuration, utilizing the Phonology Bench~\citep{suvarna2024phonologybench}. The results indicate that {\textsc BLT} achieves superior performance compared to the baseline Llama 3 1T tokenizer-based model on this specific task.

\paragraph{CUTE}
To evaluate character-level capabilities, we benchmark {\textsc BLT} using the CUTE suite~\citep{edman2024cute}, which consists of multiple tasks divided into three main groups: compositional understanding, orthographic similarity, and sequence manipulation. This evaluation is notably demanding for most models that use tokenization, as these typically retain information about token spellings yet encounter difficulties leveraging this knowledge to actually manipulate text. As shown in Table~\ref{tab:char_tasks}, {\textsc BLT}-Entropy achieves more than a 25-point lead over both BPE Llama 3 variants on this benchmark. Notably, our model displays outstanding skill in character manipulation, attaining 99.9\% accuracy on both spelling-related tasks. Such dramatic gains, especially considering that {\textsc BLT} is trained with sixteen times less data than Llama 3.1, suggest that BPE-based models find it challenging to capture character-level structures. Figure \ref{fig:cute_examples} presents several cases where Llama 3 with BPE tokenization encounters failures, while our {\textsc BLT} approach is successful. Among the evaluated tasks, only word deletion and insertion see BPE outperforming BLT; although such manipulations might pose greater challenges to models processing data at the byte level, the performance difference is modest, and it may be inherently easier for character-based models to generalize from characters up to words than the reverse. For consistency, all experiments employ identical evaluation procedures and retain the original prompts sourced from Huggingface. With further prompt optimization, BPE models might be able to improve their performance.

\paragraph{Low Resource Machine Translation} We assess {\textsc BLT} on translation tasks involving both directions across six major language families and twenty-one lower-resource languages, which encompass a variety of scripts, utilizing the FLORES-101 benchmark~\citep{goyal2022flores}. SentencePiece BLEU scores are presented in Table~\ref{tab:flores}. The evaluation reveals that {\textsc BLT} delivers superior performance relative to a system leveraging the Llama 3 tokenizer, with gains of 2 BLEU points for translations into English and 0.5 BLEU points for translations from English. While the performance of {\textsc BLT} is on par with or marginally surpasses Llama 3 in prominent language pairs, it clearly surpasses Llama 3 in many cases within lower-resource language families. These findings highlight the benefit of byte-level modeling for robust generalization over underrepresented byte sequences.

\subsection{Training {\textsc BLT} Using Llama 3 Initialization}
\label{sec:distillation}
We investigate a strategy wherein {\textsc BLT} models benefit from pre-existing pretrained tokenizer-based models to facilitate more efficient and rapid training convergence. Specifically, we initialize the global transformer parameters of {\textsc BLT} using those from a pretrained Llama 3.1 model. The training protocol involves updating the global transformer’s parameters with a learning rate set to one-tenth of that used for the local encoder and local decoder components, following the number of optimization steps recommended for Llama 3. Performance outcomes, as shown in Table~\ref{tab:distill}, are compared with a baseline {\textsc BLT}. The results demonstrate that initializing {\textsc BLT} from Llama 3.1 yields marked improvements over both the original Llama 3 and the baseline {\textsc BLT}, even when all models are trained for an equivalent number of floating-point operations. Notably, relative to our {\textsc BLT}-Entropy model—which, as documented in Table~\ref{tab:evals}, was trained on a much larger dataset comprising 1T tokens—the {\textsc BLT} initialized from Llama 3.1 continues to outperform on the MMLU benchmark. This indicates that such initialization can substantially decrease the computational resources needed for training while achieving superior results.

This approach can be interpreted as adapting tokenizer-based architectures into tokenizer-free formats, thereby transforming a pre-trained LLaMA 3.1 model into a {\textsc BLT} architecture. For a thorough evaluation, we report results for the original LLaMA 3.1 model—trained on 15T tokens—in Table \ref{tab:distill} and compare its outcomes to those obtained by the {\textsc BLT} variant built from LLaMA 3. In our experiments, the modified model demonstrates a minor reduction in performance on MMLU and HumanEval, while experiencing a more pronounced decrease on additional benchmarks. These observations indicate the necessity for further investigation to better exploit the pre-trained model’s capabilities, particularly through refinements in data composition and the adjustment of related hyperparameters.

\section{Ablations and Discussion}
\label{section:ablations}
Here, we present ablation studies that underpin the architectural selections made for {\textsc BLT}, as well as the adopted patching methodology and chosen hyper-parameters for the {\textsc BLT} model with 8B parameters, trained on the {\textsc BLT}-1T dataset.

\paragraph{Entropy Model Hyper-parameters} To evaluate how different entropy model sizes and variations in context window length impact scaling dynamics, we experiment with byte-level entropy transformer models, scaling model sizes from 1 million to 100 million parameters and adjusting context windows from 64 up to 512. We generate bpb versus training flop scaling law curves by leveraging our $400m$ and $1b$ parameter {\textsc BLT} models trained on the Llama-2 dataset, as visualized in \autoref{fig:entropy_ablation}. Our results indicate that scaling up both model size and context window generally leads to improved performance; however, benefits plateau once the model size exceeds approximately 50 million parameters.

\paragraph{Types of Patching}

We conduct an ablation study on the four patching approaches described in Section~\ref{section:patching}: (1) Strided Patching utilizing strides of 4 and 6, (2) Patching at whitespace characters, (3) Patching according to BPE tokens as defined by the Llama 3 tokenizer, and (4) Entropy-based patching that leverages a compact byte-level language model.

Although dynamic patching shortens the actual sequence length, we ensure that the context length remains consistent across all patching strategies by regulating the sequence length. During both training and inference, each model is exposed to an equivalent expected number of bytes per sequence, thereby eliminating any potential confounds related to differences in context size. The findings from these ablation studies are summarized in Figure~\ref{fig:scaling-compute-opt}. All alternative patching methods yield better performance compared to static patching, with space patching performing almost as well as dynamic entropy based patching.

\autoref{tab:evals_ablation} displays the results of benchmark assessments for {\textsc BLT} models, contrasting tokenizer-based models, space patching, and entropy-based patching approaches, all trained for the optimal number of iterations on the Llama 2 dataset~\citep{dubey2024llama}. While space patching offers a more straightforward method that bypasses the need to execute an entropy model dynamically during training, our analysis shows that the improvements identified with entropy-based patching in the context of scaling trends (see Section~\ref{section:scaling}) also extend to a range of downstream benchmark evaluations.\footnote{The space patching results shown derive from prior experiments without cross-attention; however, comparable patterns are seen when cross-attention is included.}

\paragraph{Cross-Attention} 
As presented in~\autoref{tab:crossattn_ablations_1b}, we systematically remove and vary the application of cross-attention within different layers of both the encoder and decoder modules in {\textsc BLT}. For encoder cross-attention, we evaluate three distinct strategies for initializing queries: (1) assigning an identical, learned embedding to all global states, (2) utilizing a hash-based embedding derived from the patch’s constituent bytes, and (3) applying a pooling operation over the encoder's hidden states corresponding to the patch bytes at the specified layer.

Our results indicate that incorporating cross-attention within the \textit{decoder} yields the greatest performance gains. Applying cross-attention in the encoder leads to marginal benefits, and these are observed only when queries are initialized via pooling. Moreover, cross-attention proves especially advantageous on Common-Crawl data, with the effect being more pronounced for larger patch sizes.

\paragraph{n-gram Hash Embeddings} 

We experiment with various n-gram hash embedding vocabulary sizes—specifically 0, 100K, 200K, and 400K—and summarize the outcomes in Table~\ref{tab:ngram_ablations_1b}. Our observations reveal that hash embeddings yield improvements across every evaluated domain, with the most notable effects seen for Wikipedia and Github (showing a 0.04 bpb improvement, compared to 0.01 bpb after 15k steps at the 8B scale). When increasing hash counts from 300K to 500K at the 8B scale, the resultant performance shift at 15k steps is minimal, around 0.001 bpb. These results suggest that hash embeddings are crucial for elevating the performance of {\textsc BLT} to the levels of tokenizer-based models, yet benefits start to plateau beyond 300K hashes. Furthermore, our data suggest that hash embeddings and cross-attention contribute improvements in mostly non-overlapping areas, as their advantages manifest on distinct datasets.

\paragraph{Local Model Hyperparameters} Table \ref{tab:local_layers} presents an ablation study examining different choices for the number of layers within the local encoder and decoder modules. Our results indicate that, in combination with hash n-gram embeddings, {\textsc BLT} performs effectively even when the encoder is highly minimal—consisting of only a single layer—while the decoder benefits from being comparatively deeper.

\section{Related Work}
\label{section:related_work}

\textbf{Character-Level RNNs:} Character Language Modeling has remained a focal task since the advent of neural architectures~\citep{sutskever2011generating,mikolov2012subword,graves2013generating}, largely because such models inherently handle out-of-vocabulary tokens, thereby eliminating the necessity for explicit back-off strategies. In a related effort, \cite{kim2016character} trained a neural network operating exclusively at the character level on the input, employing convolutional and highway layers to generate representations that are then processed by LSTM-based RNNs. This approach achieved performance on par with the top RNN language models in English at the time, and even surpassed them in languages with richer morphology—a frequently cited benefit of character-level LLMs. Similarly, \cite{kenter2018byte} investigated byte-level LSTM architectures for machine comprehension tasks, demonstrating superiority over word-level models, particularly for Turkish and Russian, which are morphologically complex. Extending this paradigm, \cite{zhang2015character} explored character-level convolutional models for classification problems, observing improved results over word-based approaches in certain scenarios. In further developments, \citet{chung2019hierarchical} introduced hierarchical LSTM models equipped with boundary detectors at every tier, thereby uncovering the implicit hierarchical structure in textual data and further enhancing character-level language modeling. For machine translation, ByteNet proposed by \cite{kalchbrenner2016neural} utilizes convolutional neural networks at the character level, distinguishing itself from attention-based models.

\textbf{Character-Level Transformers:} The introduction of transformer architectures employing attention mechanisms~\citep{vaswani2017attention} alongside subword tokenization methods~\citep{sennrich-etal-2016-neural} has led to marked advancements in the capabilities of neural models for language modeling and standardized evaluation tasks. Nonetheless, segmenting inputs into words or sub-word units imposes a predetermined level of abstraction for model processing. In pursuit of integrating the benefits of transformer approaches with the encouraging outcomes observed in character-level language modeling, \citet{al2019character} constructed very deep transformer networks. By utilizing auxiliary loss functions, they successfully trained transformer-based models that achieved superior performance compared to earlier LSTM-driven character language models. Despite this progress, their models still exhibited a notable performance gap relative to large language models trained at the word level. Similarly, GPT-2~\citep{radfordgpt2} reported that, when trained on expansive datasets such as the 1 Billion Word benchmark, byte-level language models lagged behind their word-level counterparts in competitiveness.

Although \cite{Choe2019BridgingTG} showed that transformer-based LLMs operating at the byte level can surpass subword-level LLMs with a similar parameter count, these byte-level models require considerably more computation and significantly longer training times. In a related effort, \cite{el2020characterbert} introduced CharFormer, a BERT variant that generates word-level representations through convolutions over character embeddings, yielding performance gains in medical domain applications but at the cost of increased computational demand. In another study, \cite{clark2022canine} proposed CANINE, an encoder-only model with 150M parameters that processes character sequences directly. CANINE incorporates a deep transformer core akin to our global model architecture, complemented by a local transformer and strided convolutions to reduce the character input length, resulting in superior performance compared to the token-level encoder-only model mBERT on various multilingual benchmarks. ByT5~\citep{xue2022byt5} investigated byte-level encoder-decoder models that eschew patching operations entirely; while this approach led to greater noise robustness and competitive performance with subword-tokenizer models using only a quarter of the data, the absence of patching necessitated attention computation across all bytes, greatly increasing computational costs. The challenge with byte-level models lies in the longer input sequences compared to subword units, making computational scaling more difficult. Most recently, leveraging the Mamba Architecture~\citep{gu2023mamba}, which enables retention of a fixed-size memory state across extensive contexts, \cite{wang2024mambabyte} trained a byte-level Mamba-based model without any patching, achieving performance that surpasses byte-level transformers in flop-matched conditions at a 350M parameter scale in terms of bits-per-byte over multiple datasets.

\textbf{Patching-based approaches:} Leveraging patching methods can significantly reduce the excessive computational costs typically associated with byte-level LLMs, without substantial loss in performance. Several studies have shown promising results for such techniques, particularly in settings with modest model sizes and training data volumes. For example, \cite{nawrot-etal-2022-hierarchical} investigate fixed (static) patching through downsampling and upsampling operations, introducing the hourglass transformer architecture. At the 150M parameter scale, their model surpasses existing byte-level baselines. Building on this, \cite{nawrot-etal-2023-efficient} present enhancements via dynamic patching strategies, incorporating both an end-to-end learnable boundary predictor, a tokenizer-supervised boundary prediction model, and an entropy-driven patching mechanism resembling {\textsc BLT}. They demonstrate that these methods can exceed the performance of contemporaneous vanilla transformers on language modeling benchmarks with models containing approximately 40M parameters trained on about 400M tokens. In a different line of inquiry, \cite{lester2024training} study the effects of using sequences compressed through arithmetic coding, which allows for greater compression than traditional BPE. By adopting an equal-info windows methodology, their models surpass byte-level baselines in language modeling, though they fall short compared to subword-based systems.

Our research is heavily influenced by MegaByte~\citep{yu2023megabyte}, a causal decoder-only large language model (LLM) that relies on fixed, static patching, where input bytes are grouped and concatenated into patches for processing. This architecture then employs a localized model within the decoder to reconstruct bytes from these patches. The original MegaByte work demonstrates that this approach achieves performance on par with tokenizer-based models when evaluated at a 1B parameter scale using a dataset comprising approximately 400B bytes. In our study, we perform ablations of MegaByte in all our experimental setups and observe that static patching underperforms relative to the best compute-efficient, tokenizer-based models in a controlled flop regime; we further illustrate how {\textsc BLT} successfully narrows this performance gap. Additionally, \cite{slagle2024spacebyte} reach similar conclusions regarding MegaByte, proposing modifications to the static patching strategy by including whitespace and space-like byte boundaries, as well as introducing a local encoder component. Their findings indicate that such enhancements yield improvements over tokenized transformer models within compute-controlled settings for particular domains such as Github and arXiv at the 1B parameter level. We expand on these results with our own experiments, showing that additional innovations in architecture are necessary to advance byte-level models to the point where they are truly competitive with state-of-the-art token-based systems like Llama 3.

\section{Limitations and Future Work}

In this study, we adopt the training step counts identified as optimal for Llama~3~\citep{dubey2024llama} when making architectural decisions. It is important to note, however, that these scaling laws were derived based on BPE-level transformer models, which may result in less-than-ideal ratios between data volume and parameter count for {\textsc BLT}. The development of specific scaling laws tailored to {\textsc BLT}, which could yield even more advantageous scaling behaviors for our model, is left for future investigation. Furthermore, a significant portion of our experimental evaluation has been restricted to models containing up to 1B parameters; thus, as we extend our work to encompass 8B parameter models and beyond, it is conceivable that the set of optimal architectural configurations may shift, potentially leading to greater performance improvements at larger scales.

Current transformer codebases and libraries are primarily optimized for tokenizer-based transformer models. Although our work includes theoretical flop-matched experiments and incorporates efficient solutions (for example, FlexAttention) to support layers that differ from the standard transformer setup, the actual wall-clock performance of our implementations may still lag behind that of tokenizer-based approaches, indicating that there is potential for additional optimization.

Whereas {\textsc BLT} relies on an independently trained entropy model for patching, integrating the training of the patching model into a unified end-to-end process presents a promising avenue for future exploration. In Section \ref{sec:distillation}, we introduce preliminary experiments suggesting positive outcomes when applying the “byte-ifying” approach to tokenizer-based architectures such as Llama 3, which have been pre-trained on over 10T tokens, by initializing the global transformer and subsequently freezing its parameters. Continued investigation in this area could lead to strategies that not only preserve the advantages of bytefying but also surpass the performance of these tokenizer-based systems—without necessitating complete retraining from the beginning.

\section{Conclusion}
\label{section:conclusion}

In this work, we introduce the Byte Latent Transformer (\textbf{BLT}), an innovative model architecture that challenges the standard reliance on predetermined vocabulary tokenization found in most large language models. {\textsc BLT} leverages a novel, adaptive mechanism that dynamically learns to aggregate bytes into variable-length patches, thereby enabling computational efforts to be distributed in proportion to the complexity of the input data. This approach results in notable gains in both computational efficiency and robustness. Our comprehensive scaling experiments show that {\textsc BLT} attains performance on par with tokenization-based systems such as Llama 3, tested at sizes up to 8B parameters and 4T bytes, and is able to reduce inference flops by as much as 50\% with only minimal compromises in evaluation benchmarks. In addition, {\textsc BLT} paves the way for a fresh scaling strategy, making it possible to simultaneously scale up model size and patch size under a constant inference cost, which is particularly beneficial for compute constraints typical of real-world deployments. By processing inputs directly at the byte level, {\textsc BLT} enhances the model's capacity to generalize to less frequent and noisier data, delivering notable robustness benefits and more nuanced handling of sub-word components. Collectively, our findings highlight {\textsc BLT} as a compelling substitute for standard tokenization-dependent methods, offering a scalable, efficient, and resilient foundation for next-generation language models.

\end{document}